%% kapa: exam=1 round=34 year=2023 period=1 form=A subject=economics subject_ko=경제학원론 seq=10 number=50
%% source: Q-Net 감정평가사 기출문제 (cst00309) artlSeq=5211789 file=raw/1008/5211789/2023년도 제34회 감정평가사 1차시험 1교시.pdf
%% conversion: kordoc 4.24.51 -> tools/gen1.py (auto); review=proofread-sonnet; flags=none
\begin{question}{50}
게임이론에 관한 설명으로 옳은 것은?
\begin{choices}
  \item 내쉬균형은 상대방의 전략에 관계없이 자신에게 가장 유리한 전략을 선택하는 것을 말한다.
  \item 복점시장에서의 내쉬균형은 하나만 존재한다.
  \item 어떤 게임에서 우월전략균형이 존재하지 않더라도 내쉬균형은 존재할 수 있다.
  \item 순차게임에서는 내쉬조건만 충족하면 완전균형이 된다.
  \item 승자의 불행(winner's curse) 현상을 방지하기 위해 최고가격입찰제(first-price sealed-bid auction)가 도입되었다.
\end{choices}
\end{question}
\answer{3}
